\section{Introdução}

As mesmas emissões de dióxido de carbono (CO\textsubscript{2}) podem produzir distribuições setoriais distintas conforme o critério de atribuição. Pela produção, são contabilizadas na atividade em que ocorrem; pela demanda, são associadas ao setor do produto final, considerando sua cadeia de produção. As emissões da geração de eletricidade, por exemplo, podem ser registradas no setor elétrico ou atribuídas aos produtos finais cuja produção requer essa eletricidade. Comparar essas perspectivas ajuda a relacionar os setores emissores às demandas que mobilizam sua produção \parencite{marques2012,millerBlair2009}.

Este trabalho pergunta: \textbf{quais relações intersetoriais compõem as diferenças entre as emissões atribuídas pela produção e pela demanda na economia brasileira de 2015, e como essas relações se organizam em termos de volume, concentração e conectividade?} O objetivo é caracterizar essa estrutura e discutir suas implicações para identificar oportunidades de coordenação da mitigação. Utiliza-se a Matriz de Insumo-Produto (MIP) nacional de 2015, com 67 atividades \parencite{ibge2015mip}, combinada com coeficientes setoriais de emissão.

Os totais setoriais não mostram, isoladamente, como as origens emissoras se relacionam aos produtos finais. Uma pegada elevada pode concentrar-se em poucas origens ou distribuir-se entre muitas; um grande emissor pode atender a demandas finais de diferentes atividades. Além disso, um setor pode apresentar uma diferença pequena entre as contas e, ao mesmo tempo, ter grandes volumes de emissões próprias atribuídas a outros produtos finais e de emissões de outras atividades atribuídas aos seus produtos finais. Identificar as contrapartes e a distribuição dos pesos permite distinguir esses perfis, além de comparar suas posições nos rankings de emissões.

A proposta articula a contabilidade ambiental à análise de redes ponderadas. \textcite{millerBlair2009} fundamentam o modelo insumo-produto e a extensão ambiental que relaciona produção e demanda final. \textcite{sanguinetAzzoni2024} examinam emissões nas cadeias produtivas regionais brasileiras sob diferentes perspectivas de atribuição. No campo das redes, \textcite{huang2019} constroem uma rede de emissões para a China e combinam filtragem de ligações, indicadores de posição e comunidades. Este trabalho adapta a representação de atribuição e a seleção de ligações ao recorte nacional brasileiro, com foco na relação entre as contas pela produção e pela demanda.

A contribuição pretendida é identificar quais contrapartes compõem as diferenças entre as duas contas e se as relações de atribuição se concentram em poucas atividades ou se distribuem entre muitas. A análise tem caráter descritivo e exploratório. Na representação em rede, as atividades são os nós e as ligações dirigidas associam origens emissoras a produtos finais, com pesos dados pelos volumes de emissões atribuídos. Serão examinados os volumes, a concentração dos pesos e a conectividade de uma rede reduzida por um critério explícito de seleção de ligações, denominada \textit{backbone}. Para essa seleção, será avaliado o filtro de disparidade de \textcite{serrano2009}. As medidas de concentração serão calculadas antes da filtragem, pois o número de ligações selecionadas não representa, por si só, a diversidade das relações originais.

O recorte abrange emissões da produção nacional associadas à demanda doméstica e às exportações, excluindo emissões estrangeiras. Os coeficientes de \textcite{sanguinetAzzoni2024} constituem a referência inicial de um cenário exploratório, com validação ambiental e monetária pendente, detalhada na seção~\ref{sec:dados}. Por isso, o trabalho ainda não sustenta rankings definitivos de emissores brasileiros.

A discussão de política se concentrará em conjuntos de setores e demandas que mereçam investigação conjunta. Intervenções nos processos produtivos e no uso de energia e materiais ao longo das cadeias são frentes complementares de mitigação discutidas por \textcite{rissman2020}. A rede poderá informar quais relações examinar, mas não estimará a eficácia ou o custo de instrumentos de política.

As seções seguintes apresentam os dados, desenvolvem a metodologia, expõem os resultados e discutem suas implicações e limitações, encerrando com as conclusões.
